% !PRINTED 247
\weberpage{247}
%<<CONT
\weq{4}{27D = 4(a^2+12c)^3 - (2a^3-72ac+27b^2)^2}
is determined, is at the same time the discriminant of the biquadratic equation (1), and the signs of the quantities $u$, $v$, $w$ are restricted by the condition
\weq{5}{uvw = -b}
[\S.\,37, (3)].

If the discriminant is negative, then the equation (3), just like (1), has two conjugate imaginary roots.

If $D$ is positive and all four roots $\alpha$, $\beta$, $\gamma$, $\delta$ are real, then $u$, $v$, $w$ also become real, and hence their squares, i.\,e. the roots of (3), are \gesperrt{positive}.

If all four roots are imaginary, say $\alpha$ conjugate to $\beta$ and $\gamma$ to $\delta$, then it follows from (2) that $v$ and $w$ are purely imaginary and $u$ is real; hence in this case the equation (3) has one positive and two negative roots.

In both cases, however, if $b$ vanishes, one of the quantities $u$, $v$, $w$ may also be equal to zero.

Having regard to the results of the last section on the cubic equation, we thus arrive at the following result:

\gesperrt{The necessary and sufficient condition for the existence of four distinct real roots is}
\weq{6}{D > 0, \quad a < 0, \quad a^2 - 4c > 0.}

In all other cases in which $D$ is positive, the equation has four imaginary roots.

We can also carry the discussion through completely for $D=0$, that is, in the case of the equality of two roots.

If we assume $\gamma = \delta$, then, since we have made the assumption $\alpha+\beta+\gamma+\delta = 0$, it follows from (2) that
\[
2v = 2w = \alpha - \beta
\]
\[
2u = \alpha+\beta-2\gamma = 2(\alpha+\beta)
\]
and from this, by (3),
\weq{7}{\begin{aligned} -4a &= 2(\alpha+\beta)^2 + (\alpha-\beta)^2\\ 16(a^2-4c) &= (\alpha-\beta)^2[8(\alpha+\beta)^2+(\alpha-\beta)^2]. \end{aligned}}
To dispose first of the case in which also $\alpha=\beta$, we obtain from (7) for this the condition $a^2-4c=0$, and the first equation (7) shows that $a$ is negative if $\alpha$ and $\gamma$
%CONT>>
